\documentclass[11pt]{amsart}

\usepackage[T1]{fontenc}
\usepackage{lmodern}
\usepackage{microtype}
\usepackage{mathtools,amssymb,amsthm}
\usepackage[hidelinks]{hyperref}

\hypersetup{
  pdftitle={A four-element counterexample to Conjecture 6.1 of Mohan and Neetu
            on two disjoint abelian sets}
}

\newtheorem{theorem}{Theorem}
\newtheorem{lemma}[theorem]{Lemma}
\theoremstyle{remark}
\newtheorem{remark}[theorem]{Remark}
\newtheorem*{conjA}{Conjecture 6.1 (Mohan and Neetu)}
\newtheorem*{exA}{Example 6.1 (Mohan and Neetu)}

\newcommand{\K}{K}

\title[A counterexample to a two-abelian-set conjecture]
{A Four-Element Counterexample to Conjecture~6.1 of Mohan and Neetu
 on Two Disjoint Abelian Sets}

\date{}

\subjclass[2020]{11P70, 20F60, 20F65, 11B13}
\keywords{small doubling, product set, torsion-free group, right-ordered group,
Klein bottle group, Baumslag--Solitar group, corrigendum}

\begin{document}

\begin{abstract}
Call a subset $S$ of a group \emph{abelian} when $\langle S\rangle$ is abelian.
Conjecture~6.1 of Mohan and Neetu (arXiv:2607.11194v1, \S6) asserts that if
$S=A\cup B$ is a finite subset of a torsion-free group, $A$ and $B$ disjoint
abelian sets, and $|S^{2}|\le3k-3$, then $\langle S\rangle$ is abelian; the
usage in that paper fixes $k=|S|$. The conjecture is false. In the Klein bottle
group $\mathrm{BS}(1,-1)$, written additively in two coordinates, the
four-element set $S=A\cup B$ with $A=\{(-1,1),(-1,3)\}$ and
$B=\{(1,1),(1,3)\}$ has $A$ and $B$ disjoint abelian sets, $|S^{2}|=9=3|S|-3$,
and $\langle S\rangle$ nonabelian. This set is not new: it is the union of the
first two of the three parts of Example~6.1 of that paper, printed two lines
below the conjecture, at parameter value~$2$. No theorem of the paper is
contradicted.
\end{abstract}

\maketitle

\section{The statement}

Throughout, $G$ is a group with identity $e$, $S\subseteq G$ is finite, and
\[
  S^{2}=\{s_{1}s_{2}:s_{1},s_{2}\in S\}
\]
is the product set over \emph{ordered} pairs. Following \cite[\S1]{MN} ---
``a set $S$ is said to be abelian if $\langle S\rangle$ is an abelian subgroup
of $G$'' --- a set is called \emph{abelian} when the subgroup it generates is
abelian. Freiman, Herzog, Longobardi and Maj proved that if $G$ is
\emph{orderable}, $|S|\ge3$ and $|S^{2}|\le3|S|-3$, then $S$ is abelian
\cite[Theorem~1.3]{FHLM}, quoted as \cite[Theorem~2.3]{MN}; the last section of
\cite{MN}, ``Future Problems'', proposes ways to keep that conclusion while
weakening the hypothesis. The first of them is:

\begin{conjA}
Let $S$ be a nonempty finite subset of a torsion-free group $G$ such that
$S=A\cup B$, where $A$ and $B$ are disjoint abelian sets. If
$|S^{2}|\le 3k-3$, $\langle S\rangle$ is abelian.
\end{conjA}

\noindent
The conjecture does not define $k$; the usage in \cite{MN} fixes $k=|S|$.
Example~6.1, printed two lines below the conjecture and quoted in full in \S2,
asserts ``$|S^{2}|=3|S|-3$'' of a set satisfying the antecedent, and
\cite[Theorem~2.9]{MN} writes ``where $k=|S|$''. This is not a formality: under
the alternative reading $k=|A|$ the set exhibited below would fall outside the
antecedent, since there $3|A|-3=3<9=|S^{2}|$.

\section{The witness}

Let $\K=\mathbb{Z}\rtimes_{-1}\mathbb{Z}=\mathrm{BS}(1,-1)$ be the Klein bottle
group, in the coordinates used by \cite[\S5]{MN}: elements are pairs
$(r,n)\in\mathbb{Z}^{2}$ and
\begin{equation}\label{eq:law}
  (r,n)(s,m)=\bigl(r+(-1)^{n}s,\;n+m\bigr),
  \qquad
  (r,n)^{-1}=\bigl(-(-1)^{n}r,\,-n\bigr),
\end{equation}
with $e=(0,0)$. This is the group and the operation of \cite{MN}:
\cite[\S5]{MN} writes $\mathrm{BS}(1,q)\cong\mathbb{Z}[1/q]\rtimes\mathbb{Z}$
with $(r,n)(s,m)=(r+q^{n}s,n+m)$, and the examples of its \S6 are stated in a
group the authors write $\mathbb{Z}\rtimes\mathbb{Z}$; that $q=-1$ there is
fixed by their Example~6.2, which records the identity $(i,1)(j,1)=(i-j,2)$,
and for $q=1$ the group would be abelian. All coordinates therefore stay in
$\mathbb{Z}$ and every computation below is on pairs of integers.

\begin{lemma}\label{lem:tf}
$\K$ is torsion-free and nonabelian.
\end{lemma}

\begin{proof}
By \eqref{eq:law} the second coordinate is additive, so $(r,n)^{j}$ has second
coordinate $jn$: if $n\ne0$ and $j\ne0$ this is nonzero, so
$(r,n)^{j}\ne e$. If $n=0$ then $(r,0)^{j}=(jr,0)\ne e$ whenever $r\ne0$ and
$j\ne0$. Hence no element other than $e$ has finite order. Finally
$(-1,1)(1,1)=(-2,2)$ while $(1,1)(-1,1)=(2,2)$.
\end{proof}

\begin{lemma}\label{lem:pc}
If the elements of a set $X\subseteq G$ commute pairwise, then $\langle X\rangle$
is abelian, i.e.\ $X$ is an abelian set in the sense of \cite{MN}.
\end{lemma}

\begin{proof}
The centraliser $C_{G}(X)$ is a subgroup containing $X$, hence containing
$\langle X\rangle$; so $\langle X\rangle\subseteq C_{G}(\langle X\rangle)$.
\end{proof}

Two lines below Conjecture~6.1, the authors print:

\begin{exA}
Let
\begin{align*}
  A&=\{(-1,1),(-1,3),\ldots,(-1,2k-1)\},\\
  B&=\{(1,1),(1,3),\ldots,(1,2k-1)\},\\
  C&=\{(-1,2),(1,2),(-1,4),(1,4),\ldots,(-1,2k),(1,2k)\}
\end{align*}
be subsets of the group $\mathbb{Z}\rtimes\mathbb{Z}$. Let $S=A\cup B\cup C$.
Then, we have $|S^{2}|=3|S|-3$. It is easy to observe that $\langle A\rangle$,
$\langle B\rangle$, and $\langle C\rangle$ are abelian, whereas $\langle
S\rangle$ is nonabelian.
\end{exA}

\noindent
Set the parameter the authors call $k$ there --- not the $k$ of the conjecture
--- to $2$, and delete the third part $C$. What is left is the following.

\begin{theorem}\label{thm:main}
In $\K$ let
\[
  A=\{(-1,1),(-1,3)\},\qquad B=\{(1,1),(1,3)\},\qquad S=A\cup B .
\]
Then $A$ and $B$ are disjoint abelian sets, $|S|=4$,
\[
  S^{2}=\bigl\{(d,2j):d\in\{-2,0,2\},\ j\in\{1,2,3\}\bigr\},
  \qquad |S^{2}|=9=3|S|-3,
\]
and $\langle S\rangle$ is nonabelian. Since $\K$ is torsion-free by
Lemma~\ref{lem:tf} and $9\le3\cdot4-3$, Conjecture~6.1 of \cite{MN} is false.
\end{theorem}

\begin{proof}
The four listed pairs are distinct, and $A\cap B=\varnothing$ because the first
coordinates $-1$ and $1$ differ. Every second coordinate occurring in $S$ is
odd, so $(-1)^{n}=-1$ throughout and \eqref{eq:law} gives, for
$\varepsilon,\delta\in\{-1,1\}$ and odd $n,n'$,
\[
  (\varepsilon,n)(\delta,n')=(\varepsilon-\delta,\,n+n'),
  \qquad \varepsilon-\delta\in\{-2,0,2\}.
\]
Taking $\delta=\varepsilon$ gives $(0,n+n')$, which is symmetric in $n,n'$; so
the elements of $A$ commute pairwise and so do those of $B$, and by
Lemma~\ref{lem:pc} both are abelian sets.

For the product set: with $n,n'\in\{1,3\}$ the sum $n+n'$ takes exactly the
values $2,4,6$, and $S$ contains both $(-1,n)$ and $(1,n)$ for $n=1$ and
$n=3$, so at each of those three levels all four choices of
$(\varepsilon,\delta)$ occur and $\varepsilon-\delta$ realises each of
$-2,0,2$. That is the displayed set, and it has $3\cdot3=9$ elements, while
$3|S|-3=3\cdot4-3=9$.

Finally $(-1,1)(1,1)=(-2,2)$ and $(1,1)(-1,1)=(2,2)$, and both factors lie in
$S$, so $\langle S\rangle$ is nonabelian.
\end{proof}

\noindent
No minimality is claimed: nothing here decides whether a witness of
cardinality less than four exists, nor whether $3|S|-3$ is the least doubling
attainable by such a configuration, nor anything about host groups other than
$\K$.

\begin{remark}[the contrast drawn after the conjecture]
In \cite{MN} the conjecture is followed by ``Whereas this may not hold for
three abelian sets. Consider the following example.'', introducing Example~6.1.
Read literally, the first sentence asserts only that the conclusion may fail
for three abelian sets, which it does; but the contrast it draws with the
two-set case is misleading, since by Theorem~\ref{thm:main} the conclusion
fails for two abelian sets as well, and the witness comes from the very example
that sentence introduces.
\end{remark}

\begin{remark}[no theorem of \cite{MN} is contradicted]
The set of Theorem~\ref{thm:main} lies outside the hypotheses of the two
positive results of \cite{MN} nearest to it. \cite[Theorem~2.10]{MN} requires,
among other hypotheses, $|A|\ge4$ and $|B|\ge3$, whereas here $|A|=|B|=2$;
and \cite[Theorem~2.7]{MN} requires
$|(A\cup B)^{2}|\le3|A\cup B|-4=8$, whereas $|S^{2}|=9$.
\end{remark}

\section{Verification}

Every object used above is printed above, and the decisive computation is the
sixteen ordered products of the four-element set $S$, carried out in the proof
of Theorem~\ref{thm:main} by the identity
$(\varepsilon,n)(\delta,n')=(\varepsilon-\delta,n+n')$ for odd $n,n'$; a reader
needs no code.

The program \texttt{verify.py} accompanies this note and its recorded output is
\texttt{verify.output.txt}. It re-implements \eqref{eq:law} in exact integer
arithmetic --- the sign $(-1)^{n}$ taken as a parity lookup, so that no float
can appear --- forms the sixteen products of $S$ one by one, compares $S^{2}$
with the nine pairs displayed in Theorem~\ref{thm:main} \emph{as a set of
coordinate pairs, not merely by cardinality}, and confirms $|S^{2}|=9=3|S|-3$,
that $A$ and $B$ are disjoint commuting pairs, and that
$(-1,1)(1,1)\ne(1,1)(-1,1)$. Torsion-freeness of $\K$ is confirmed there only
on a bounded box of coordinates and exponents; the unbounded statement is
Lemma~\ref{lem:tf}. The recorded output also contains checks on objects not
used above, among them the three-set Example~6.1 of \cite{MN} and the family of
which the set of Theorem~\ref{thm:main} is one member; nothing above depends on
them.

\begin{thebibliography}{9}

\bibitem{FHLM}
G.~A. Freiman, M.~Herzog, P.~Longobardi, and M.~Maj,
\emph{Small doubling in ordered groups},
J. Aust. Math. Soc. \textbf{96} (2014), 316--325.
Its Theorem~1.3 is quoted as \cite[Theorem~2.3]{MN}, which is the form used
here.

\bibitem{MN}
Mohan and Neetu,
\emph{On small doubling in right-ordered groups and Baumslag--Solitar
groups~--~II},
arXiv:2607.11194v1, 13 July 2026.
\href{https://arxiv.org/abs/2607.11194}{arXiv:2607.11194}.
Statement numbers are those the source compiles to.

\end{thebibliography}

\end{document}
